\section{Related Work}
\label{sec:related_work}

\textbf{Ontology-based IR and food ontologies.}
Classical ontology-based retrieval enriches document and query representations with explicit semantic structure. Vallet et al.~\cite{vallet2005ontology} combine keyword search with ontology annotations, Castells et al.~\cite{castells2007adaptation} adapt the vector-space model to ontology-based ranking, and Fern\'{a}ndez et al.~\cite{fernandez2011semantic} extend this line to larger-scale settings, collectively showing that concept-level representations reduce vocabulary mismatch and improve retrieval precision. In the food domain, FoodOn~\cite{dooley2018foodon} provides a large harmonized food ontology for traceability and data integration, while FoodKG~\cite{haussmann2019foodkg} links recipes, ingredients, and nutrients in a knowledge graph used for explainable recommendation. Chen et al.~\cite{chen2021personalized} further formulate personalized food recommendation as constrained question answering over a food knowledge graph. However, these food-oriented studies target recommendation or question answering rather than ranking-based retrieval, and the classical ontology-IR studies predate dense retrieval. Neither line integrates ontology with modern neural retrieval pipelines, particularly for low-resource languages such as Vietnamese.

\textbf{Food retrieval under neural frameworks.}
Recent work treats food search as a retrieval problem within neural pipelines. Recipe-MPR~\cite{zhang2023recipempr} benchmarks multi-aspect preference-based recipe retrieval, highlighting the difficulty of queries that span multiple attributes. Hu et al.~\cite{hu2024bridging} propose a cross-cultural recipe retrieval framework that benefits from query reformulation and reranking. KERL~\cite{mohbat2025kerl} integrates food knowledge graphs with large language models for personalized recipe recommendation, confirming that structured food knowledge remains useful in LLM-based systems. Meanwhile, dense retrieval pipelines~\cite{karpukhin2020dense, gao2023rag} have become the dominant paradigm for knowledge-intensive tasks, yet analyses of retrieval failure modes~\cite{barnett2024seven, edge2024local} show that dense retrieval struggles with class-level and compositional queries that require structured reasoning. None of these studies places ontology at the core of the retrieval pipeline as a pre-embedding enrichment mechanism for food-specific queries.

\textbf{Ingredient substitution and dish similarity.}
Ingredient substitution has been studied through both knowledge-graph and learning-based approaches. Shirai et al.~\cite{shirai2021identifying} identify substitutions using a food knowledge graph, while Fatemi et al.~\cite{fatemi2023learning} learn substitution patterns from recipe corpora. For dish similarity, Park et al.~\cite{park2021flavorgraph} construct FlavorGraph, a large-scale food-chemical graph that captures flavor compatibility. Sun et al.~\cite{sun2023think} and He et al.~\cite{he2024g} demonstrate more broadly that coupling retrieval with knowledge-graph traversal improves reasoning on compositional queries. However, existing substitution methods lack constraint-aware filtering through a class hierarchy, and dish similarity is typically computed as flat Jaccard over ingredient bags without distinguishing within-class from cross-class differences.

\textbf{Position of the present work.}
Our work bridges these three lines by: (1)~constructing a Vietnamese-specific ontology from a 10K-dish dataset, (2)~integrating it into a dense-retrieval pipeline as a pre-embedding enrichment mechanism for three food retrieval tasks, and (3)~providing an explicit ablation that isolates the contribution of hierarchy, typed relations, and inference rules.
